\documentclass[11pt, a4paper]{article}

% bfw-style.tex — gemeinsame Style-Definitionen für alle Handouts/Aufgabenhefte
% Voraussetzung: Kompilation mit xelatex (wegen fontspec)

\usepackage[ngerman]{babel}
\usepackage{geometry}
\geometry{a4paper, top=2.2cm, bottom=2.2cm, left=2.3cm, right=2.3cm}

\usepackage{fontspec}
% Fallback-Kette: Source Sans Pro -> Calibri -> Segoe UI -> Latin Modern Sans
% (Latin Modern ist immer mit MiKTeX vorhanden)
\IfFontExistsTF{Source Sans Pro}{%
  \setmainfont{Source Sans Pro}[Ligatures=TeX]%
}{\IfFontExistsTF{Calibri}{%
  \setmainfont{Calibri}[Ligatures=TeX]%
}{\IfFontExistsTF{Segoe UI}{%
  \setmainfont{Segoe UI}[Ligatures=TeX]%
}{%
  \setmainfont{Latin Modern Sans}[Ligatures=TeX]%
}}}
\IfFontExistsTF{Source Code Pro}{%
  \setmonofont{Source Code Pro}[Scale=0.92]%
}{\IfFontExistsTF{Consolas}{%
  \setmonofont{Consolas}[Scale=0.92]%
}{%
  \setmonofont{Latin Modern Mono}[Scale=0.92]%
}}

\usepackage{xcolor}
% Farbpalette: hell, freundlich, modern
\definecolor{bfwPrimary}{HTML}{0EA5A4}    % Petrol/Türkis - Primär
\definecolor{bfwPrimaryDark}{HTML}{0F766E} % dunkler Petrol für Headings
\definecolor{bfwAccent}{HTML}{F59E0B}     % Amber - Akzent
\definecolor{bfwSuccess}{HTML}{10B981}    % Grün - Erfolg / Lösung
\definecolor{bfwWarn}{HTML}{F97316}       % Orange - Achtung
\definecolor{bfwInk}{HTML}{0F172A}        % fast Schwarz
\definecolor{bfwInkSoft}{HTML}{475569}    % gedämpftes Grau
\definecolor{bfwBgSoft}{HTML}{F8FAFC}     % off-white
\definecolor{bfwBgTeal}{HTML}{ECFEFF}     % helles Türkis
\definecolor{bfwBgAmber}{HTML}{FEF3C7}    % helles Gelb
\definecolor{bfwBgRose}{HTML}{FFE4E6}     % helles Rosé für Eselsbrücken

\usepackage[colorlinks=true, linkcolor=bfwPrimaryDark, urlcolor=bfwPrimaryDark]{hyperref}
\usepackage{enumitem}
\setlist[itemize]{leftmargin=*, itemsep=2pt, topsep=2pt}
\setlist[enumerate]{leftmargin=*, itemsep=2pt, topsep=2pt}

\usepackage[most]{tcolorbox}
\usepackage{booktabs}
\usepackage{tikz}
\usetikzlibrary{decorations.pathmorphing, positioning, shapes.geometric, arrows.meta}

% Merkkästchen — wichtige Definition / Kernpunkt
\newtcolorbox{merksatz}[1][]{%
  enhanced, breakable,
  colback=bfwBgTeal,
  colframe=bfwPrimary,
  boxrule=0.6pt,
  arc=4pt,
  left=10pt, right=10pt, top=8pt, bottom=8pt,
  fonttitle=\bfseries\color{bfwPrimaryDark},
  title={\faLightbulb\ Merksatz},
  #1
}

% Eselsbrücke — Mnemoniks
\newtcolorbox{eselsbruecke}[1][]{%
  enhanced, breakable,
  colback=bfwBgRose,
  colframe=bfwAccent,
  boxrule=0.6pt,
  arc=4pt,
  left=10pt, right=10pt, top=8pt, bottom=8pt,
  fonttitle=\bfseries\color{bfwWarn},
  title={Eselsbrücke},
  #1
}

% Beispiel-Box
\newtcolorbox{beispiel}[1][]{%
  enhanced, breakable,
  colback=bfwBgSoft,
  colframe=bfwInkSoft,
  boxrule=0.4pt,
  arc=3pt,
  left=10pt, right=10pt, top=8pt, bottom=8pt,
  fonttitle=\bfseries\color{bfwInk},
  title={Beispiel},
  #1
}

% Lösungs-Box (für Aufgabenhefte)
\newtcolorbox{loesung}[1][]{%
  enhanced, breakable,
  colback=bfwBgAmber!40,
  colframe=bfwSuccess,
  boxrule=0.6pt,
  arc=3pt,
  left=10pt, right=10pt, top=8pt, bottom=8pt,
  fonttitle=\bfseries\color{bfwSuccess!70!black},
  title={Lösung},
  #1
}

% Glossar-Eintrag
\newcommand{\glossareintrag}[2]{%
  \par\noindent\textbf{\color{bfwPrimaryDark}#1} \enspace #2\par\smallskip
}

% Section-Stil — etwas farbiger als Standard
\usepackage{titlesec}
\titleformat{\section}
  {\Large\bfseries\color{bfwPrimaryDark}}
  {\thesection}{0.6em}{}
\titleformat{\subsection}
  {\large\bfseries\color{bfwPrimary}}
  {\thesubsection}{0.5em}{}
\titleformat{\subsubsection}
  {\normalsize\bfseries\color{bfwInk}}
  {\thesubsubsection}{0.4em}{}

% TikZ-Stil "skizzenhaft" — etwas weicher als Rough.js, aber stabil
\tikzset{
  roughbox/.style = {
    draw=bfwPrimaryDark, thick, fill=bfwBgTeal,
    rounded corners=4pt, inner sep=6pt
  },
  roughaccent/.style = {
    draw=bfwAccent, thick, fill=bfwBgAmber,
    rounded corners=4pt, inner sep=6pt
  },
  roughline/.style = {
    draw=bfwInkSoft, thick, -{Stealth[length=2.5mm]}
  }
}

% Bullet-Symbole (bewusst kein FontAwesome — vermeidet Font-Installations-Probleme)
\newcommand{\faLightbulb}{$\bullet$}
\newcommand{\faBookmark}{$\bullet$}

% Header/Footer
\usepackage{fancyhdr}
\pagestyle{fancy}
\fancyhf{}
\renewcommand{\headrulewidth}{0pt}
\fancyfoot[L]{\small\color{bfwInkSoft}\thepage}
\fancyfoot[R]{\small\color{bfwInkSoft} BFW M-064 Beschaffung im Büromanagement}


\title{\vspace*{-1.5cm}\Huge\bfseries\color{bfwPrimaryDark}Aufgabenheft Tag~3\\[0.4em]
  \large\color{bfwInkSoft}\textnormal{Bestellungen und Kaufverträge --- Übungen im IHK-Klausurformat}}
\author{\textsf{Beschaffung im Büromanagement}\\
\textsf{\small\copyright\ Can Siebert\,\textperiodcentered\,2026}}
\date{Selbstlernmaterial \textperiodcentered{} 13.05.2026}

\begin{document}
\maketitle
\thispagestyle{empty}

\begin{merksatz}[title={\bfseries So nutzen Sie dieses Aufgabenheft}]
Bearbeiten Sie alle Aufgaben zunächst \emph{ohne} in die Lösungen zu schauen.
Beachten Sie die \emph{Operatoren} und schreiben Sie Antworten in vollständigen Sätzen.
Lösungen ab Seite~\pageref{loesungen}.
\end{merksatz}

\tableofcontents
\vspace{1em}
\hrule

\section{Aufgaben}

\subsection*{Aufgabe~1 --- Geschäftsfähigkeit (10 Punkte)}

\noindent\textbf{\color{bfwAccent}YouTube-Vorbereitung:}\enspace \href{https://www.youtube.com/watch?v=0DNydbv2Qb8}{Beschränkte Geschäftsfähigkeit Minderjähriger} (\url{https://www.youtube.com/watch?v=0DNydbv2Qb8})

\medskip
\begin{enumerate}[label=(\alph*)]
  \item \textbf{Nennen} Sie die drei Stufen der Geschäftsfähigkeit nach BGB. \hfill (3~P.)
  \item \textbf{Erläutern} Sie den ,,Taschengeldparagraphen'' (§\,110 BGB). \hfill (3~P.)
  \item \textbf{Beurteilen} Sie folgende Fälle und begründen Sie: \hfill (4~P.)
    \begin{enumerate}[label=\arabic*., leftmargin=*]
      \item Ein 6-jähriges Kind kauft einen Kaugummi für 50~Cent vom Taschengeld.
      \item Ein 16-jähriger Auszubildender bestellt online ein Tablet für 400~€.
    \end{enumerate}
\end{enumerate}

\subsection*{Aufgabe~2 --- Rechtsgeschäfte und Trennungsprinzip (12 Punkte)}

\noindent\textbf{\color{bfwAccent}YouTube-Vorbereitung:}\enspace \href{https://www.youtube.com/watch?v=8huvlRCYCT0}{Rechtsgeschäfte einfach erklärt} (\url{https://www.youtube.com/watch?v=8huvlRCYCT0})

\medskip
\begin{enumerate}[label=(\alph*)]
  \item \textbf{Unterscheiden} Sie einseitige und mehrseitige Rechtsgeschäfte und nennen Sie pro Art ein Beispiel. \hfill (4~P.)
  \item \textbf{Erläutern} Sie an einem Kaufvertrag den Unterschied zwischen Verpflichtungs\-geschäft und Verfügungs\-geschäft. \hfill (5~P.)
  \item \textbf{Beurteilen} Sie: Warum macht das Trennungs\-prinzip im Alltag Sinn? \hfill (3~P.)
\end{enumerate}

\subsection*{Aufgabe~3 --- Vertragsfreiheit und ihre Grenzen (10 Punkte)}

\noindent\textbf{\color{bfwAccent}YouTube-Vorbereitung:}\enspace \href{https://www.youtube.com/watch?v=gPManz0Ke1o}{Was ist ein Vertrag? Angebot und Annahme} (\url{https://www.youtube.com/watch?v=gPManz0Ke1o})

\medskip
\begin{enumerate}[label=(\alph*)]
  \item \textbf{Nennen} Sie vier Aspekte der Vertragsfreiheit. \hfill (4~P.)
  \item \textbf{Erläutern} Sie die Schranken nach §\,134 BGB (Verbotsgesetz) und §\,138 BGB (Sittenwidrigkeit). \hfill (4~P.)
  \item \textbf{Beurteilen} Sie: Ist ein mündlich abgeschlossener Kaufvertrag rechtlich gültig? \hfill (2~P.)
\end{enumerate}

\subsection*{Aufgabe~4 --- Zustandekommen eines Kaufvertrags (12 Punkte)}

\noindent\textbf{\color{bfwAccent}YouTube-Vorbereitung:}\enspace \href{https://www.youtube.com/watch?v=g7j4NkgtiNo}{Kaufvertrag: Zustandekommen einfach erklärt} (\url{https://www.youtube.com/watch?v=g7j4NkgtiNo})

\medskip
Frau Meier sieht im Onlineshop von \emph{Bürofix AG} Aktenordner für 4{,}90~€ pro Stück
und bestellt 50~Stück.

\begin{enumerate}[label=(\alph*)]
  \item \textbf{Erläutern} Sie, durch welche zwei Willenserklärungen ein Kaufvertrag zustande kommt. \hfill (3~P.)
  \item \textbf{Prüfen} Sie: Ist die Anzeige der Aktenordner im Onlineshop bereits ein bindendes Angebot oder eine \emph{invitatio ad offerendum}? Begründen Sie. \hfill (4~P.)
  \item \textbf{Beschreiben} Sie, ab welchem Schritt der Kaufvertrag tatsächlich zustande kommt. \hfill (3~P.)
  \item \textbf{Nennen} Sie die zwei \emph{essentialia negotii} eines Kaufvertrags. \hfill (2~P.)
\end{enumerate}

\subsection*{Aufgabe~5 --- AGB-Einbeziehung (8 Punkte)}

\noindent\textbf{\color{bfwAccent}YouTube-Vorbereitung:}\enspace \href{https://www.youtube.com/watch?v=lueRR4lHgbY}{AGB einfach erklärt für Azubis} (\url{https://www.youtube.com/watch?v=lueRR4lHgbY})

\medskip
\begin{enumerate}[label=(\alph*)]
  \item \textbf{Nennen} Sie die drei Voraussetzungen, unter denen AGB wirksam in einen Vertrag einbezogen werden (§\,305 Abs.\,2 BGB). \hfill (3~P.)
  \item \textbf{Erläutern} Sie, was eine ,,überraschende Klausel'' ist und was passiert, wenn eine solche in den AGB steht. \hfill (3~P.)
  \item \textbf{Beurteilen} Sie: Sind individuell ausgehandelte Vertragsbedingungen AGB-Recht unterworfen? \hfill (2~P.)
\end{enumerate}

\subsection*{Aufgabe~6 --- AGB-Inhaltskontrolle (Praxisfall) (8 Punkte)}

\noindent\textbf{\color{bfwAccent}YouTube-Vorbereitung:}\enspace \href{https://www.youtube.com/watch?v=UucDMI8z0OY}{§\,305 BGB: Allgemeine Geschäftsbedingungen} (\url{https://www.youtube.com/watch?v=UucDMI8z0OY})

\medskip
In den AGB der \emph{Bürofix AG} steht: \emph{,,Reklamationen müssen innerhalb von
24~Stunden nach Lieferung schriftlich angezeigt werden, sonst gilt die Lieferung als
angenommen.''}

\begin{enumerate}[label=(\alph*)]
  \item \textbf{Beurteilen} Sie diese Klausel im Hinblick auf §\,307 BGB (unangemessene Benachteiligung). \hfill (4~P.)
  \item \textbf{Erläutern} Sie, welche gesetzliche Regelung statt der unwirksamen AGB-Klausel greift. \hfill (2~P.)
  \item \textbf{Begründen} Sie, warum die AGB-Inhaltskontrolle besonders für Verbraucher (B2C) wichtig ist. \hfill (2~P.)
\end{enumerate}

\vspace{1em}
\noindent\textbf{Punkte gesamt: 60}

\newpage

\section{Lösungen}\label{loesungen}

\subsection*{Lösung Aufgabe~1}
\begin{loesung}
\textbf{(a)} Drei Stufen:
\begin{itemize}[itemsep=0pt]
  \item \emph{Geschäftsunfähig}: Kinder unter 7~Jahren (§\,104 BGB), dauerhaft Geistesgestörte. Willenserklärungen nichtig (§\,105 BGB).
  \item \emph{Beschränkt geschäftsfähig}: Minderjährige zwischen 7 und 17~Jahren (§\,106 BGB). Nur lediglich rechtlich vorteilhafte Geschäfte selbständig wirksam.
  \item \emph{Voll geschäftsfähig}: Personen ab dem 18.~Geburtstag.
\end{itemize}

\textbf{(b)} Der Taschengeldparagraph (§\,110 BGB) regelt: Ein Vertrag eines Minderjährigen
ohne Eltern­zustimmung ist trotzdem von Anfang an wirksam, wenn die vertragsmäßige
Leistung mit Mitteln bewirkt wird, die ihm zur freien Verfügung überlassen wurden
(z.\,B.\ Taschengeld). Das soll Minderjährige am üblichen Einkauf teilnehmen lassen,
ohne ständige Eltern­zustimmung.

\textbf{(c)}
\begin{enumerate}[itemsep=0pt]
  \item Das 6-jährige Kind ist \emph{geschäftsunfähig} (§\,104 BGB). Sein Kaufvertrag ist nichtig (§\,105 BGB) --- unabhängig vom geringen Wert. Es hat einen Rückgabeanspruch.
  \item Der 16-jährige ist \emph{beschränkt geschäftsfähig}. Der Vertrag über 400~€ ist schwebend unwirksam, bis die Eltern zustimmen (§\,108 Abs.\,1 BGB). Ohne Zustimmung wird er endgültig unwirksam. Ausnahme: §\,110 BGB, wenn die 400~€ aus seinem freiverfügbaren Geld stammen.
\end{enumerate}
\end{loesung}

\subsection*{Lösung Aufgabe~2}
\begin{loesung}
\textbf{(a)}
\begin{itemize}[itemsep=0pt]
  \item \emph{Einseitiges Rechtsgeschäft}: nur eine Willenserklärung. Beispiele: Kündigung, Testament, Anfechtung.
  \item \emph{Mehrseitiges Rechtsgeschäft (Vertrag)}: zwei oder mehr übereinstimmende Willenserklärungen. Beispiele: Kaufvertrag, Mietvertrag, Arbeitsvertrag.
\end{itemize}

\textbf{(b)} Beim Kaufvertrag passieren rechtlich zwei getrennte Vorgänge (Trennungsprinzip):
\begin{itemize}[itemsep=0pt]
  \item \emph{Verpflichtungsgeschäft}: Der Kaufvertrag selbst (§\,433 BGB). Er begründet die Pflichten --- Verkäufer schuldet Übergabe + Eigentum, Käufer schuldet Geld + Abnahme.
  \item \emph{Verfügungsgeschäft}: Die Erfüllung der Pflichten. Übergabe + Eigentumsübertragung der Sache (§\,929 BGB) auf der einen, Geldzahlung auf der anderen Seite.
\end{itemize}

\textbf{(c)} Das Trennungsprinzip macht im Alltag Sinn, weil Vertragsschluss und Erfüllung
zeitlich auseinanderfallen können. Wenn ich heute einen Aktenordner kaufe und morgen
geliefert bekomme, ist der Vertrag heute geschlossen, aber das Eigentum geht erst morgen
über. Außerdem schützt das Abstraktionsprinzip vor Folgekosten: Wenn der Vertrag
nichtig ist, war die Eigentumsübertragung trotzdem wirksam --- der Verkäufer muss
zwar das Eigentum zurückübertragen, aber das ist juristisch ein eigenständiger
Vorgang.
\end{loesung}

\subsection*{Lösung Aufgabe~3}
\begin{loesung}
\textbf{(a)} Vier Aspekte der Vertragsfreiheit:
\begin{itemize}[itemsep=0pt]
  \item Abschlussfreiheit (ob und mit wem Verträge geschlossen werden)
  \item Inhaltsfreiheit (was im Vertrag steht)
  \item Formfreiheit (mündlich, schriftlich, Handschlag)
  \item Aufhebungsfreiheit (einvernehmliche Auflösung)
\end{itemize}

\textbf{(b)}
\begin{itemize}[itemsep=0pt]
  \item \emph{§\,134 BGB --- Verbotsgesetz}: Verträge gegen ein gesetzliches Verbot sind nichtig (z.\,B.\ Drogen- oder Schwarzarbeitsvertrag).
  \item \emph{§\,138 BGB --- Sittenwidrigkeit/Wucher}: Verträge gegen die guten Sitten sind nichtig. Wucher (Spezialfall): Ausnutzen einer Notlage zu auffälligem Missverhältnis von Leistung und Gegenleistung (z.\,B.\ Kreditvertrag mit 200\,\% Zinsen p.\,a.).
\end{itemize}

\textbf{(c)} Ja --- ein mündlicher Kaufvertrag ist grundsätzlich rechtlich gültig
(Formfreiheit). Ausnahmen sind gesetzliche Form­vorschriften, z.\,B.\ Grundstücks­kauf
nach §\,311b BGB (notariell). Im Alltagsgeschäft (Bäcker, Tankstelle) sind mündliche
Verträge die Regel.
\end{loesung}

\subsection*{Lösung Aufgabe~4}
\begin{loesung}
\textbf{(a)} Ein Kaufvertrag kommt zustande durch zwei übereinstimmende Willens­
erklärungen: \textbf{Antrag} (Angebot, §\,145 BGB) und \textbf{Annahme} (rechtzeitig
und vorbehaltlos). Beide müssen sich auf denselben Vertragsgegenstand beziehen.

\textbf{(b)} Die Anzeige im Onlineshop ist eine \emph{invitatio ad offerendum} ---
nur eine Aufforderung zur Abgabe eines Angebots, kein bindender Antrag. Begründung:
Der Shop muss sich gegen unbeschränkte Bestellungen schützen (begrenzter Lagerbestand,
Bonität des Käufers, Kapazität). Würde die Schaufensterauslage / Online-Anzeige als
Antrag gelten, müsste der Shop jede Annahme akzeptieren.

\textbf{(c)} Der Vertragsschluss erfolgt in folgenden Schritten:
\begin{enumerate}[itemsep=0pt]
  \item Frau Meier sendet die Bestellung ab $\to$ macht damit den \textbf{Antrag}.
  \item Eine automatische ,,Eingangsbestätigung'' ist meist noch nicht die Annahme.
  \item Die \textbf{Annahme} erfolgt mit der ausdrücklichen Bestellbestätigung
    (Versand­bestätigung) des Shops --- jetzt ist der Vertrag perfekt.
\end{enumerate}

\textbf{(d)} Die zwei (von drei) essentialia negotii sind \textbf{Ware} (welcher
Aktenordner) und \textbf{Preis} (4{,}90~€/Stück). Daneben braucht man die Vertragspartner
--- damit drei Hauptpunkte gesamt.
\end{loesung}

\subsection*{Lösung Aufgabe~5}
\begin{loesung}
\textbf{(a)} Drei Voraussetzungen nach §\,305 Abs.\,2 BGB:
\begin{enumerate}[itemsep=0pt]
  \item \emph{Hinweis} auf die AGB durch den Verwender (vor oder bei Vertragsschluss).
  \item Möglichkeit zur \emph{zumutbaren Kenntnisnahme} (Aushang, Link, Beilage).
  \item \emph{Einverständnis} der anderen Partei (ausdrücklich oder konkludent).
\end{enumerate}

\textbf{(b)} Eine \emph{überraschende Klausel} (§\,305c BGB) ist eine AGB-Bestimmung,
die nach den Umständen so ungewöhnlich ist, dass die andere Partei nicht damit rechnen
musste --- z.\,B.\ versteckte Klauseln, untypische Inhalte, irreführende Überschriften.
\emph{Folge:} Sie wird \textbf{nicht} Vertragsbestandteil. Der Vertrag bleibt im Übrigen
wirksam, die überraschende Klausel fällt einfach weg.

\textbf{(c)} Nein. Individuell ausgehandelte Vertragsbedingungen sind nach §\,305 Abs.\,1
Satz~3 BGB \emph{keine} AGB. Wenn beide Parteien tatsächlich verhandelt haben (nicht nur
,,vorgelegt und akzeptiert''), unterliegt die Klausel nicht dem AGB-Recht --- nur dem
allgemeinen Vertragsrecht und §\,138 BGB (Sittenwidrigkeit) als Schranke.
\end{loesung}

\subsection*{Lösung Aufgabe~6}
\begin{loesung}
\textbf{(a)} Die Klausel verkürzt die gesetzlichen Mängelrechte (Verjährung in der Regel
2~Jahre nach §\,438 BGB) auf 24~Stunden --- das ist eine drastische Verkürzung. Nach
§\,307 BGB liegt eine \textbf{unangemessene Benachteiligung} der anderen Partei vor:
24~Stunden reichen oft nicht einmal, um die Ware ordnungsgemäß zu prüfen, geschweige
denn um eine schriftliche Reklamation zu formulieren. Die Klausel widerspricht zudem
dem Grundgedanken des Mängelrechts. \emph{Ergebnis:} Die Klausel ist nach §\,307 BGB
\textbf{unwirksam}.

\textbf{(b)} An ihre Stelle tritt das gesetzliche Recht. Konkret gilt dann: §\,438 BGB
(Verjährung der Mängelansprüche, 2~Jahre bei beweglichen Sachen) und §\,377 HGB
(unverzügliche Mängelrüge bei beidseitigen Handelskäufen --- ,,unverzüglich'' wird in
der Praxis mit einigen Werktagen interpretiert).

\textbf{(c)} Die AGB-Inhaltskontrolle ist besonders wichtig für Verbraucher (B2C), weil
sie typischerweise die schwächere Partei sind: Sie haben kein Verhandlungsgewicht, lesen
AGB selten und können einzelne Klauseln nicht ändern. Ohne Inhaltskontrolle könnten
Unternehmen die schwächere Position der Verbraucher systematisch ausnutzen. Im B2B
zwischen Kaufleuten nimmt das Gesetz mehr Selbstverantwortung an --- aber auch dort
gibt es Grenzen.
\end{loesung}

\end{document}
